\begingroup
% Kiểu cột L cho tabularx (khai báo trong nhóm riêng, theo quy ước của báo cáo).
\newcolumntype{L}[1]{>{\hsize=#1\hsize\raggedright\arraybackslash}X}

\section{Video minh chứng và bản đồ phân cảnh}
\label{sec:video}

Ngoài phần thuyết minh và các ảnh minh chứng, nhóm nộp kèm một video thực hành
được dựng từ các clip quay màn hình và quay trực tiếp thiết bị (thư mục lưu trữ
clip gốc và bản dựng cuối nằm trong thư mục làm việc của nhóm). Video dài khoảng
40 phút, chia thành 26 cảnh theo 7 phần, mỗi cảnh là một clip riêng và đều có
thuyết minh của thành viên trực tiếp thao tác, đáp ứng yêu cầu ``video có thuyết
minh và thể hiện người cấu hình''. Nội dung video bám sát hai mô hình đã trình
bày: mô hình hạ tầng trên router thật (Phần 1--4) và mô hình Ad-hoc trên máy ảo
(Phần 5).

\begin{table}[htbp]
  \centering\footnotesize
  \caption{Bản đồ phân cảnh của video minh chứng và liên hệ với báo cáo.}
  \label{tab:bandoscanh}
  \begin{tabularx}{\linewidth}{@{} L{0.55} L{0.95} L{1.8} L{0.7} @{}}
    \toprule
    \textbf{Phần} & \textbf{Cảnh} & \textbf{Nội dung và điều được chứng minh} &
    \textbf{Mục báo cáo} \\
    \midrule
    0 -- Mở đầu & 01 & Giới thiệu nhóm, mục tiêu bài tập và phạm vi demo &
      Mục giới thiệu \\
    1 -- Giới thiệu bộ định tuyến & 02--03 & Giới thiệu mẫu thiết bị TP-Link Archer C20 v5 cài
      OpenWrt, mặt trước/mặt sau, cách nối cáp quang, cổng WAN và các cổng LAN;
      máy bàn cắm dây vào cổng LAN & \ref{subsec:thiet-bi-thuc-te} \\
    2 -- Cấu hình cơ bản & 04--08 & Đăng nhập LuCI lần đầu và đặt mật khẩu quản
      trị; đặt hostname \texttt{N09-Router} và múi giờ; đặt địa chỉ
      \ac{lan}/DHCP; cấu hình WAN; bật Wi-Fi hai băng tần, đặt \ac{ssid} và bảo
      mật WPA2/CCMP & \ref{sec:configuration-guide} \\
    3 -- Kết nối thiết bị (mô hình hạ tầng) & 09--14 & Máy bàn chuyển sang Wi-Fi,
      đọc IP/cổng/kênh/bảo mật, chạy \texttt{nmcli} và ping; laptop Windows tắt
      MAC ngẫu nhiên rồi kết nối; Android kết nối băng 5~GHz; kiểm tra trên LuCI
      bằng \emph{Associated Stations} và \emph{Active DHCP Leases}; ping chéo
      giữa các thiết bị & \ref{subsec:evidence-infra},
      \ref{tab:ping-crosscheck} \\
    4 -- Chức năng nâng cao & 15--21 & Ẩn ESSID; tạo mạng khách tách riêng dải
      \texttt{192.168.20.0/24}; lọc MAC; cấp IP tĩnh theo MAC; chặn tên miền qua
      DNS; giới hạn giờ truy cập bằng luật tường lửa; chuyển tiếp cổng 8080, NAT
      (Masquerading), traceroute và sao lưu cấu hình &
      \ref{subsec:cauhinh-nangcao} \\
    5 -- Mô hình Ad-hoc & 22--25 & Giải thích vì sao phải dùng máy ảo (Windows,
      Android, macOS không hỗ trợ \ac{ibss}); tạo hai card Wi-Fi ảo bằng
      \texttt{mac80211\_hwsim}; tách vào hai network namespace; đưa hai nút vào
      \ac{ibss} \texttt{N09-AdHoc} ở kênh~6; kiểm tra \texttt{iw info},
      \texttt{station dump} và ping hai chiều &
      \ref{sec:adhoc-config}, \ref{subsec:evidence-adhoc} \\
    6 -- Kết & 26 & Tổng kết những gì đã cấu hình và thử được, so sánh ngắn hai
      mô hình & \ref{sec:arch} \\
    \bottomrule
  \end{tabularx}
\end{table}

Video cho thấy ba điều mà ảnh tĩnh không thể hiện được: (i) trạng thái router
trước và sau mỗi lần áp dụng cấu hình, kể cả cảnh báo ``router chưa có mật
khẩu'' ở lần đăng nhập đầu tiên; (ii) phản ứng tức thời của thiết bị đầu cuối
khi bị chặn (mất mạng khách, bị lọc MAC, bị chặn theo giờ); và (iii) toàn bộ
quá trình dựng mạng Ad-hoc từ lúc chưa có card mạng nào đến khi hai nút ping được
nhau. Các số liệu được đọc trực tiếp trên màn hình trong video cũng chính là số
liệu dùng trong Bảng~\ref{tab:dhcp-leases},
Bảng~\ref{tab:ping-crosscheck} và Bảng~\ref{tab:adhoc-results}.

\endgroup
